\chapter{Overview}
\label{chap:overview}

This blueprint records the formalization of a simplified proof of the quantum
soundness of the low individual degree test of Ji, Natarajan, Vidick, Wright,
and Yuen \cite{Ji2020LowIndividualDegree}.  The simplified proof keeps the
test, the strategies, and the three consistency conclusions of the original
paper, but it removes the sampling parameter $k$ from the final theorem,
rounds measurements to projective measurements with a linear loss following
de~la~Salle \cite{deLaSalle2022Orthogonalization}, and replaces the pasting
construction of the original paper by a first-success interpolation whose
analysis uses a single positive contraction.  The quantitative bound is
\[
	21000\,m^4\max\{1,d^2\}\bigl(\eps^{1/64}+(d/q)^{1/64}\bigr).
\]
The statement of the original paper, with its sampling parameter, is recovered
as a corollary (Corollary~\ref{cor:main-formal-sampling}).

An $m$-variate polynomial $g\colon\F_q^m\to\F_q$ has individual degree $d$ if
each variable occurs with degree at most $d$ in every monomial.  In the low
individual degree test, a verifier asks one prover for the value of a function
at a random point and the other prover for the restriction of the function to
a random axis-parallel line through that point, and accepts if the two answers
agree.  Classical soundness of such tests is recorded in the following two
theorems, which are external inputs to the literature and are not used by the
formalized proof.

\begin{theorem}[Raz--Safra]\label{thm:raz-safra}
	\uses{def:polyfunc}
	Suppose deterministic classical Provers~$\mathrm{A}$ and~$\mathrm{B}$ pass
	the $k=2$ surface-versus-point low-degree test with probability $1-\eps$.
	Then there exists a degree-$d$ polynomial $g\colon\F_q^m\to\F_q$ such that
	\[
		\Prb_{u\sim\F_q^m}[g(u)=a]\ge1-\eps-\poly(m)\cdot\poly(d/q),
	\]
	where $a$ denotes the point answer of Prover~$\mathrm{A}$.
\end{theorem}
\begin{proof}
	See Raz and Safra \cite{RazSafra1997}.
\end{proof}

\begin{theorem}[Polishchuk--Spielman]\label{thm:classical-test-soundness}
	\uses{def:low-individual-degree-test}
	Suppose deterministic classical Provers~$\mathrm{A}$ and~$\mathrm{B}$ pass
	the low individual degree test with probability $1-\eps$.  Then there
	exists a polynomial $g\colon\F_q^m\to\F_q$ with individual degree $d$ such
	that
	\[
		\Prb_{u\sim\F_q^m}[g(u)=a]
		\ge1-\poly(m)\cdot\bigl(\poly(\eps)+\poly(d/q)\bigr),
	\]
	where $a$ denotes the point answer of Prover~$\mathrm{A}$.
\end{theorem}
\begin{proof}
	See Polishchuk and Spielman \cite{PolishchukSpielman1994}; the
	multivariate statement follows from their bivariate analysis.
\end{proof}

\begin{remark}[A degree-$(d+1)$ strategy that passes with probability $1-\frac1m$]\label{ex:bad-individual-degree-example}
	Assume $d+1<q$ and consider $h(x_1,\ldots,x_m)=x_1^{d+1}$.  On a point
	question $u$, Prover~$\mathrm{A}$ answers $h(u)$.  On an axis-parallel line
	$\ell=\{u+xe_i:x\in\F_q\}$, Prover~$\mathrm{B}$ answers $h|_\ell$ when
	$i\ne1$, since $h$ is then constant along $\ell$, and answers the zero
	polynomial when $i=1$.  The verifier accepts whenever $i\ne1$, so the
	strategy passes with probability at least $1-\eps$ for $\eps=\frac1m$.
	After fixing $x_2,\ldots,x_m$, the difference between $h$ and any
	polynomial $g$ of individual degree $d$ is a nonzero univariate polynomial
	of degree $d+1$, so by Schwartz--Zippel
	\[
		\Prb_{u\sim\F_q^m}[g(u)=a]\le1-m\eps+\frac{d+1}q.
	\]
	Thus a dependence on $m$ is unavoidable in the soundness error.
\end{remark}

\begin{theorem}[Main theorem, informal]\label{thm:main-informal}
	\uses{def:low-individual-degree-test, def:projective-strategy, def:polymeasurements}
	Suppose Provers~$\mathrm{A}$ and~$\mathrm{B}$ pass the two-prover degree-$d$
	low individual degree test with probability $1-\eps$ using a
	finite-dimensional projective strategy, and let $A=\{A^u_a\}$ be the point
	measurements.  Then there is a projective measurement $G=\{G_g\}$ whose
	outcomes are polynomials of individual degree $d$ such that
	\[
		\E_{u\sim\F_q^m}\sum_{a\in\F_q}\sum_{g:g(u)=a}
		\bra\psi A^u_a\ot G_g\ket\psi
		\ge1-\poly(m,d)\cdot\bigl(\poly(\eps)+\poly(d/q)\bigr).
	\]
\end{theorem}
\begin{proof}
	\uses{thm:main-formal}
	This is the informal form of Theorem~\ref{thm:main-formal}.
\end{proof}

The proof is an induction on the dimension $m$.  Each step applies
self-improvement to the measurements produced for the parallel slices,
dilates the resulting sub-measurements simultaneously to projective
sub-measurements on a common enlarged space, pastes them by interpolation,
and compresses the pasted measurement back to the original space
(Chapter~\ref{chap:induction}).  The final theorem symmetrizes a general
strategy, applies the induction, and rounds the resulting measurements to
projective ones (Chapter~\ref{chap:projective}).
